\documentclass[12pt,a4paper]{exam}
\usepackage[utf8]{inputenc}
\usepackage{amsmath,amssymb,amsfonts}
\usepackage{geometry}
\usepackage{enumitem}
\usepackage{graphicx}
\usepackage{hyperref}
\usepackage{xcolor}
\geometry{a4paper, top=2cm, bottom=2cm, left=2.5cm, right=2.5cm}
\hypersetup{colorlinks=true, linkcolor=blue, urlcolor=blue}
\renewcommand{\thequestion}{\textbf{\arabic{question}}}
\renewcommand{\questionlabel}{\thequestion.}
\pointname{ marks}
\pointsdroppedatright
\bracketedpoints
\setlength{\parindent}{0pt}
\setlength{\emergencystretch}{3em}
\allowdisplaybreaks
\newsavebox{\schmathbox}
\newcommand{\fitmath}[1]{%
  \sbox{\schmathbox}{$\displaystyle #1$}%
  \ifdim\wd\schmathbox>\linewidth
    \resizebox{\linewidth}{!}{\usebox{\schmathbox}}%
  \else
    \usebox{\schmathbox}%
  \fi}

\begin{document}

\thispagestyle{empty}
\begin{center}
  \textbf{\LARGE Diag Solutions} \\[0.3cm]
  \rule{\textwidth}{0.5pt}
\end{center}

\vspace{0.3cm}

\begin{questions}
\question \textbf{1.} The value of $(1 + \tan^2 \theta) \cos^2 \theta$ is

\textbf{Answer:}
(A)

\textbf{Solution:}
\begin{enumerate}
  \item Use identity $1 + \tan^2 \theta = \sec^2 \theta$
  \item Multiply $\sec^2 \theta \cdot \cos^2 \theta$
\end{enumerate}
\textbf{Derivation:}
\begin{center}
\fitmath{\sec^2 \theta \cdot \cos^2 \theta = \frac{1}{\cos^2 \theta} \cdot \cos^2 \theta = 1}
\end{center}
\hfill \textit{Introduction to Trigonometry — Trigonometric Identities}
\vspace{0.3cm}

\question \textbf{2.} If $\sin A = \cos A$, then the value of $A$ is

\textbf{Answer:}
(B)

\textbf{Solution:}
\begin{enumerate}
  \item Set $\sin A = \cos A$
  \item Divide by $\cos A$ to get $\tan A = 1$
  \item Find angle where tangent is 1
\end{enumerate}
\textbf{Derivation:}
\begin{center}
\fitmath{\begin{aligned}
\tan A = 1 \\
\implies A = 45^\circ
\end{aligned}}
\end{center}
\hfill \textit{Introduction to Trigonometry — Trigonometric Ratios of Specific Angles}
\vspace{0.3cm}

\question \textbf{3.} The value of $\text{cosec}^2 \theta - 1$ is equal to

\textbf{Answer:}
(D)

\textbf{Solution:}
\begin{enumerate}
  \item Use the identity $1 + \cot^2 \theta = \text{cosec}^2 \theta$
  \item Rearrange to solve for $\text{cosec}^2 \theta - 1$
\end{enumerate}
\textbf{Derivation:}
\begin{center}
\fitmath{\text{cosec}^2 \theta - 1 = \cot^2 \theta}
\end{center}
\hfill \textit{Introduction to Trigonometry — Trigonometric Identities}
\vspace{0.3cm}

\question \textbf{4.} Assertion (A): For an angle $\theta$, $\sin^2 \theta + \cos^2 \theta = 1$. Reason (R): This identity holds true for all values of $\theta$ where $0^\circ \le \theta \le 90^\circ$.

\textbf{Answer:}
(A)

\textbf{Solution:}
\begin{enumerate}
  \item The Pythagorean identity is a fundamental trigonometric identity.
  \item It is valid for all values of $\theta$ in the given domain.
\end{enumerate}
\textbf{Derivation:}
\begin{center}
\fitmath{\sin^2 \theta + \cos^2 \theta = 1 \text{ for all } \theta}
\end{center}
\hfill \textit{Introduction to Trigonometry — Trigonometric Identities}
\vspace{0.3cm}

\question \textbf{5.} Assertion (A): The value of $\tan 45^\circ$ is $1$. Reason (R): $\tan \theta = \frac{\sin \theta}{\cos \theta}$ and $\sin 45^\circ = \cos 45^\circ = \frac{1}{\sqrt{2}}$.

\textbf{Answer:}
(A)

\textbf{Solution:}
\begin{enumerate}
  \item Recall the trigonometric table values.
  \item Use the quotient identity $\tan \theta = \frac{\sin \theta}{\cos \theta}$ to verify.
\end{enumerate}
\textbf{Derivation:}
\begin{center}
\fitmath{\tan 45^\circ = \frac{\sin 45^\circ}{\cos 45^\circ} = \frac{1/\sqrt{2}}{1/\sqrt{2}} = 1}
\end{center}
\hfill \textit{Introduction to Trigonometry — Trigonometric Ratios of Specific Angles}
\vspace{0.3cm}

\question \textbf{6.} Assertion (A): The value of $\sin \theta$ can be $1.5$ for some angle $\theta$. Reason (R): The range of $\sin \theta$ is $[-1, 1]$.

\textbf{Answer:}
(D)

\textbf{Solution:}
\begin{enumerate}
  \item The sine function value is bounded between $-1$ and $1$.
  \item Since $1.5 > 1$, the assertion is false.
\end{enumerate}
\textbf{Derivation:}
\begin{center}
\fitmath{-1 \le \sin \theta \le 1}
\end{center}
\hfill \textit{Introduction to Trigonometry — Introduction to Trigonometry}
\vspace{0.3cm}

\question \textbf{7.} Assertion (A): $\sec^2 \theta - \tan^2 \theta = 1$. Reason (R): $1 + \tan^2 \theta = \sec^2 \theta$.

\textbf{Answer:}
(A)

\textbf{Solution:}
\begin{enumerate}
  \item The identity $1 + \tan^2 \theta = \sec^2 \theta$ is a standard NCERT identity.
  \item Rearranging this gives $\sec^2 \theta - \tan^2 \theta = 1$.
\end{enumerate}
\textbf{Derivation:}
\begin{center}
\fitmath{\begin{aligned}
1 + \tan^2 \theta = \sec^2 \theta \\
\implies \sec^2 \theta - \tan^2 \theta = 1
\end{aligned}}
\end{center}
\hfill \textit{Introduction to Trigonometry — Trigonometric Identities}
\vspace{0.3cm}

\question \textbf{8.} Assertion (A): If $\cos A = \frac{1}{2}$, then $\sin A = \frac{\sqrt{3}}{2}$. Reason (R): $\sin^2 A + \cos^2 A = 1$.

\textbf{Answer:}
(A)

\textbf{Solution:}
\begin{enumerate}
  \item Use the identity $\sin^2 A = 1 - \cos^2 A$.
  \item Substitute $\cos A = 1/2$.
  \item Calculate $\sin A = \sqrt{1 - 1/4} = \sqrt{3/4} = \sqrt{3}/2$.
\end{enumerate}
\textbf{Derivation:}
\begin{center}
\fitmath{\begin{aligned}
\sin^2 A = 1 - (\frac{1}{2})^2 = 1 - \frac{1}{4} = \frac{3}{4} \\
\implies \sin A = \frac{\sqrt{3}}{2}
\end{aligned}}
\end{center}
\hfill \textit{Introduction to Trigonometry — Trigonometric Identities}
\vspace{0.3cm}

\question \textbf{9.} If $\sin A = \cos A$, where $0 < A < 90^\circ$, find the value of $A$.

\textbf{Answer:}
\begin{center}
\fitmath{45^\circ}
\end{center}

\textbf{Solution:}
\begin{enumerate}
  \item Divide both sides by $\cos A$ to get $\tan A = 1$.
  \item Identify the angle where $\tan A = 1$ in the given range.
\end{enumerate}
\textbf{Derivation:}
\begin{center}
\fitmath{\begin{aligned}
\frac{\sin A}{\cos A} = 1 \\
\implies \tan A = 1 \\
\implies A = \tan^{-1}(1) = 45^\circ
\end{aligned}}
\end{center}
\hfill \textit{Introduction to Trigonometry — Trigonometric Ratios}
\vspace{0.3cm}

\question \textbf{10.} Evaluate $\frac{\tan 65^\circ}{\cot 25^\circ}$.

\textbf{Answer:}
\begin{center}
\fitmath{1}
\end{center}

\textbf{Solution:}
\begin{enumerate}
  \item Use the complementary angle identity $\cot \theta = \tan(90^\circ - \theta)$.
  \item Substitute $\cot 25^\circ = \tan(90^\circ - 25^\circ) = \tan 65^\circ$ and simplify.
\end{enumerate}
\textbf{Derivation:}
\begin{center}
\fitmath{\frac{\tan 65^\circ}{\tan(90^\circ - 25^\circ)} = \frac{\tan 65^\circ}{\tan 65^\circ} = 1}
\end{center}
\hfill \textit{Introduction to Trigonometry — Trigonometric Ratios of Complementary Angles}
\vspace{0.3cm}

\question \textbf{11.} Given $15 \cot A = 8$, find the value of $\sin A$.

\textbf{Answer:}
\begin{center}
\fitmath{\frac{15}{17}}
\end{center}

\textbf{Solution:}
\begin{enumerate}
  \item Find $\cot A = \frac{8}{15}$. Using the ratio $\frac{\text{adjacent}}{\text{opposite}}$, the hypotenuse is $\sqrt{8^2 + 15^2}$.
  \item Calculate hypotenuse $= \sqrt{64 + 225} = 17$, then $\sin A = \frac{\text{opposite}}{\text{hypotenuse}} = \frac{15}{17}$.
\end{enumerate}
\textbf{Derivation:}
\begin{center}
\fitmath{\begin{aligned}
\cot A = \frac{8}{15} \\
\implies \text{hyp} = \sqrt{8^2 + 15^2} = 17 \\
\implies \sin A = \frac{15}{17}
\end{aligned}}
\end{center}
\hfill \textit{Introduction to Trigonometry — Trigonometric Ratios}
\vspace{0.3cm}

\question \textbf{12.} Find the value of $2 \tan^2 45^\circ + \cos^2 30^\circ - \sin^2 60^\circ$.

\textbf{Answer:}
\begin{center}
\fitmath{2}
\end{center}

\textbf{Solution:}
\begin{enumerate}
  \item Substitute the standard values: $\tan 45^\circ = 1$, $\cos 30^\circ = \frac{\sqrt{3}}{2}$, $\sin 60^\circ = \frac{\sqrt{3}}{2}$.
  \item Simplify the expression $2(1)^2 + (\frac{\sqrt{3}}{2})^2 - (\frac{\sqrt{3}}{2})^2$.
\end{enumerate}
\textbf{Derivation:}
\begin{center}
\fitmath{2(1)^2 + \frac{3}{4} - \frac{3}{4} = 2 + 0 = 2}
\end{center}
\hfill \textit{Introduction to Trigonometry — Trigonometric Ratios of Specific Angles}
\vspace{0.3cm}

\question \textbf{13.} Prove that $\sec^2 \theta - 1 = \tan^2 \theta$.

\textbf{Answer:}
\begin{center}
\fitmath{\text{Proof based on identity} 1 + \tan^2 \theta = \sec^2 \theta}
\end{center}

\textbf{Solution:}
\begin{enumerate}
  \item Start with the fundamental identity $1 + \tan^2 \theta = \sec^2 \theta$.
  \item Subtract 1 from both sides to isolate $\tan^2 \theta$.
\end{enumerate}
\textbf{Derivation:}
\begin{center}
\fitmath{\begin{aligned}
1 + \tan^2 \theta = \sec^2 \theta \\
\implies \tan^2 \theta = \sec^2 \theta - 1
\end{aligned}}
\end{center}
\hfill \textit{Introduction to Trigonometry — Trigonometric Identities}
\vspace{0.3cm}

\question \textbf{14.} If $\sqrt{3} \tan \theta = 1$, find the value of $\sin^2 \theta - \cos^2 \theta$.

\textbf{Answer:}
\begin{center}
\fitmath{-\frac{1}{2}}
\end{center}

\textbf{Solution:}
\begin{enumerate}
  \item Find the value of $\tan \theta$ from the given equation.
  \item Determine the angle $\theta$ or use the ratio to find $\sin \theta$ and $\cos \theta$.
  \item Substitute the values into the expression $\sin^2 \theta - \cos^2 \theta$ and simplify.
\end{enumerate}
\textbf{Derivation:}
\begin{center}
\fitmath{\begin{aligned}
\tan \theta = \frac{1}{\sqrt{3}} \\
\implies \theta = 30^\circ. \\
\sin 30^\circ = \frac{1}{2}, \cos 30^\circ = \frac{\sqrt{3}}{2}. \\
\sin^2 30^\circ - \cos^2 30^\circ = (\frac{1}{2})^2 - (\frac{\sqrt{3}}{2})^2 = \frac{1}{4} - \frac{3}{4} = -\frac{2}{4} = -\frac{1}{2}.
\end{aligned}}
\end{center}
\hfill \textit{Introduction to Trigonometry — Trigonometric Ratios}
\vspace{0.3cm}

\question \textbf{15.} Prove that $\frac{\cos A}{1 + \sin A} + \frac{1 + \sin A}{\cos A} = 2 \sec A$.

\textbf{Answer:}
\begin{center}
\fitmath{\text{Proved}}
\end{center}

\textbf{Solution:}
\begin{enumerate}
  \item Take the common denominator $(1 + \sin A)\cos A$.
  \item Expand the numerator using algebraic identity $(a+b)^2 = a^2 + 2ab + b^2$.
  \item Use the identity $\sin^2 A + \cos^2 A = 1$ to simplify the expression and arrive at the result.
\end{enumerate}
\textbf{Derivation:}
\begin{center}
\fitmath{\begin{aligned}
LHS = \frac{\cos^2 A + (1 + \sin A)^2}{(1 + \sin A)\cos A} = \frac{\cos^2 A + 1 + 2\sin A + \sin^2 A}{(1 + \sin A)\cos A} \\
= \frac{(\cos^2 A + \sin^2 A) + 1 + 2\sin A}{(1 + \sin A)\cos A} = \frac{1 + 1 + 2\sin A}{(1 + \sin A)\cos A} \\
= \frac{2 + 2\sin A}{(1 + \sin A)\cos A} = \frac{2(1 + \sin A)}{(1 + \sin A)\cos A} = \frac{2}{\cos A} = 2 \sec A = RHS.
\end{aligned}}
\end{center}
\hfill \textit{Introduction to Trigonometry — Trigonometric Identities}
\vspace{0.3cm}

\question \textbf{16.} If $\tan A = \frac{3}{4}$, find all other trigonometric ratios for angle $A$.

\textbf{Answer:}
\begin{center}
\fitmath{\sin A = \frac{3}{5}, \cos A = \frac{4}{5}, \csc A = \frac{5}{3}, \sec A = \frac{5}{4}, \cot A = \frac{4}{3}}
\end{center}

\textbf{Solution:}
\begin{enumerate}
  \item Consider a right triangle where $\tan A = \frac{\text{Opposite}}{\text{Adjacent}} = \frac{3k}{4k}$.
  \item Use the Pythagorean theorem to find the hypotenuse.
  \item Calculate the remaining ratios using the sides of the triangle.
\end{enumerate}
\textbf{Derivation:}
\begin{center}
\fitmath{\begin{aligned}
\text{Let } P = 3, B = 4. \text{ By Pythagoras, } H = \sqrt{3^2 + 4^2} = \sqrt{25} = 5. \\
\sin A = \frac{P}{H} = \frac{3}{5}, \cos A = \frac{B}{H} = \frac{4}{5}, \csc A = \frac{1}{\sin A} = \frac{5}{3}, \sec A = \frac{1}{\cos A} = \frac{5}{4}, \cot A = \frac{1}{\tan A} = \frac{4}{3}.
\end{aligned}}
\end{center}
\hfill \textit{Introduction to Trigonometry — Trigonometric Ratios}
\vspace{0.3cm}

\question \textbf{17.} Prove that $\frac{\tan \theta}{1 - \cot \theta} + \frac{\cot \theta}{1 - \tan \theta} = 1 + \sec \theta \csc \theta$.

\textbf{Answer:}
\begin{center}
\fitmath{1 + \sec \theta \csc \theta}
\end{center}

\textbf{Solution:}
\begin{enumerate}
  \item Convert all terms into terms of $\sin \theta$ and $\cos \theta$.
  \item Simplify the fractions by finding a common denominator for the terms.
  \item Use the identity $\sin^2 \theta + \cos^2 \theta = 1$.
  \item Factor the numerator using the difference of cubes formula $a^3 - b^3 = (a-b)(a^2 + ab + b^2)$.
  \item Cancel the common factors and simplify the expression to the final form.
\end{enumerate}
\textbf{Derivation:}
\begin{center}
\fitmath{LHS = \frac{\frac{\sin \theta}{\cos \theta}}{1 - \frac{\cos \theta}{\sin \theta}} + \frac{\frac{\cos \theta}{\sin \theta}}{1 - \frac{\sin \theta}{\cos \theta}} = \frac{\sin^2 \theta}{\cos \theta(\sin \theta - \cos \theta)} + \frac{\cos^2 \theta}{\sin \theta(\cos \theta - \sin \theta)} = \frac{\sin^3 \theta - \cos^3 \theta}{\sin \theta \cos \theta (\sin \theta - \cos \theta)} = \frac{(\sin \theta - \cos \theta)(\sin^2 \theta + \cos^2 \theta + \sin \theta \cos \theta)}{\sin \theta \cos \theta (\sin \theta - \cos \theta)} = \frac{1 + \sin \theta \cos \theta}{\sin \theta \cos \theta} = \frac{1}{\sin \theta \cos \theta} + 1 = \csc \theta \sec \theta + 1 = RHS}
\end{center}
\hfill \textit{Introduction to Trigonometry — Trigonometric Identities}
\vspace{0.3cm}

\question \textbf{18.} If $\sec \theta + \tan \theta = p$, prove that $\sin \theta = \frac{p^2 - 1}{p^2 + 1}$.

\textbf{Answer:}
\begin{center}
\fitmath{\sin \theta = \frac{p^2 - 1}{p^2 + 1}}
\end{center}

\textbf{Solution:}
\begin{enumerate}
  \item Use the identity $\sec^2 \theta - \tan^2 \theta = 1$ to find $\sec \theta - \tan \theta$.
  \item Since $(\sec \theta + \tan \theta)(\sec \theta - \tan \theta) = 1$, then $\sec \theta - \tan \theta = 1/p$.
  \item Solve the system of two linear equations to find $\sec \theta$ and $\tan \theta$ in terms of $p$.
  \item Calculate $p^2 - 1$ and $p^2 + 1$ using the expressions found.
  \item Divide $(p^2 - 1)$ by $(p^2 + 1)$ and simplify to show it equals $\sin \theta$.
\end{enumerate}
\textbf{Derivation:}
Given  p = $\sec \theta + \tan \theta$ . We know  $\sec^2 \theta - \tan^2 \theta = 1 \implies (\sec \theta - \tan \theta) = \frac{1}{p}$ . Adding equations:  2$\sec \theta = p + \frac{1}{p} = \frac{p^2+1}{p} \implies \cos \theta = \frac{2p}{p^2+1}$ . Subtracting equations:  2$\tan \theta = p - \frac{1}{p} = \frac{p^2-1}{p} \implies \tan \theta = \frac{p^2-1}{2p}$ . Since  $\sin \theta = \tan \theta \cdot \cos \theta = \frac{p^2-1}{2p} \cdot \frac{2p}{p^2+1} = \frac{p^2-1}{p^2+1}$ .
\hfill \textit{Introduction to Trigonometry — Trigonometric Ratios}
\vspace{0.3cm}

\question \textbf{19.} Find the $11^{th}$ term of an Arithmetic Progression (AP) whose first term is $-3$ and the common difference is $4$.

\textbf{Answer:}
\begin{center}
\fitmath{37}
\end{center}

\textbf{Solution:}
\begin{enumerate}
  \item Identify the given parameters: first term $a = -3$ and common difference $d = 4$.
  \item Use the formula for the $n^{th}$ term: $a_n = a + (n-1)d$ and substitute $n=11$.
\end{enumerate}
\textbf{Derivation:}
\begin{center}
\fitmath{a_{11} = a + (11-1)d = -3 + 10(4) = -3 + 40 = 37}
\end{center}
\hfill \textit{Arithmetic Progressions — Nth term of an AP}
\vspace{0.3cm}

\question \textbf{20.} Which term of the AP $21, 18, 15, \dots$ is $-81$?

\textbf{Answer:}
\begin{center}
\fitmath{35^{th} \text{term}}
\end{center}

\textbf{Solution:}
\begin{enumerate}
  \item Identify $a = 21$ and $d = 18 - 21 = -3$.
  \item Set $a_n = -81$ and solve for $n$ using $a_n = a + (n-1)d$.
\end{enumerate}
\textbf{Derivation:}
\begin{center}
\fitmath{\begin{aligned}
-81 = 21 + (n-1)(-3) \\
\implies -102 = (n-1)(-3) \\
\implies n-1 = 34 \\
\implies n = 35
\end{aligned}}
\end{center}
\hfill \textit{Arithmetic Progressions — Nth term of an AP}
\vspace{0.3cm}

\question \textbf{21.} Find the sum of the first $10$ multiples of $5$.

\textbf{Answer:}
\begin{center}
\fitmath{275}
\end{center}

\textbf{Solution:}
\begin{enumerate}
  \item The AP is $5, 10, 15, \dots$ where $a = 5, d = 5, n = 10$.
  \item Use the sum formula $S_n = \frac{n}{2}[2a + (n-1)d]$.
\end{enumerate}
\textbf{Derivation:}
\begin{center}
\fitmath{S_{10} = \frac{10}{2}[2(5) + (10-1)5] = 5[10 + 45] = 5[55] = 275}
\end{center}
\hfill \textit{Arithmetic Progressions — Sum of first n terms}
\vspace{0.3cm}

\question \textbf{22.} If the $n^{th}$ term of an AP is given by $a_n = 3 + 2n$, find its common difference.

\textbf{Answer:}
\begin{center}
\fitmath{2}
\end{center}

\textbf{Solution:}
\begin{enumerate}
  \item Find the first term $a_1$ by substituting $n=1$.
  \item Find the second term $a_2$ by substituting $n=2$, then calculate $d = a_2 - a_1$.
\end{enumerate}
\textbf{Derivation:}
\begin{center}
\fitmath{a_1 = 3 + 2(1) = 5; \quad a_2 = 3 + 2(2) = 7; \quad d = a_2 - a_1 = 7 - 5 = 2}
\end{center}
\hfill \textit{Arithmetic Progressions — General term of an AP}
\vspace{0.3cm}

\question \textbf{23.} Find the value of $k$ for which the terms $k-1, k+3, 3k-1$ are in Arithmetic Progression.

\textbf{Answer:}
\begin{center}
\fitmath{k = 5}
\end{center}

\textbf{Solution:}
\begin{enumerate}
  \item Use the property that if $a, b, c$ are in AP, then $2b = a + c$.
  \item Substitute the terms and solve the linear equation for $k$.
\end{enumerate}
\textbf{Derivation:}
\begin{center}
\fitmath{\begin{aligned}
2(k+3) = (k-1) + (3k-1) \\
\implies 2k + 6 = 4k - 2 \\
\implies 8 = 2k \\
\implies k = 5
\end{aligned}}
\end{center}
\hfill \textit{Arithmetic Progressions — Properties of AP}
\vspace{0.3cm}

\question \textbf{24.} Find the sum of all two-digit numbers which, when divided by 4, leave 1 as a remainder.

\textbf{Answer:}
\begin{center}
\fitmath{1225}
\end{center}

\textbf{Solution:}
\begin{enumerate}
  \item Identify the sequence: The smallest two-digit number leaving remainder 1 when divided by 4 is 13 (4*3+1), and the largest is 97 (4*24+1).
  \item The sequence is 13, 17, 21, ..., 97, which is an AP with a=13, d=4, and l=97.
  \item Find the number of terms n using the formula l = a + (n-1)d.
  \item Calculate the sum using S\_n = n/2 * (a + l).
\end{enumerate}
\textbf{Derivation:}
\begin{center}
\fitmath{\begin{aligned}
97 = 13 + (n-1)4 \\
\implies 84 = (n-1)4 \\
\implies n-1 = 21 \\
\implies n = 22. S_{22} = \frac{22}{2}(13 + 97) = 11 \times 110 = 1225.
\end{aligned}}
\end{center}
\hfill \textit{Arithmetic Progressions — Sum of n terms}
\vspace{0.3cm}

\question \textbf{25.} If the sum of the first $n$ terms of an AP is $S_n = 3n^2 + 5n$, find the 16th term of the AP.

\textbf{Answer:}
\begin{center}
\fitmath{98}
\end{center}

\textbf{Solution:}
\begin{enumerate}
  \item Use the relation between the nth term and the sum of terms: $a_n = S_n - S_{n-1}$.
  \item Calculate $S_{16}$ by substituting $n=16$ into the given formula.
  \item Calculate $S_{15}$ by substituting $n=15$ into the given formula.
  \item Subtract $S_{15}$ from $S_{16}$ to find $a_{16}$.
\end{enumerate}
\textbf{Derivation:}
\begin{center}
\fitmath{a_{16} = S_{16} - S_{15} = [3(16)^2 + 5(16)] - [3(15)^2 + 5(15)] = [3(256) + 80] - [3(225) + 75] = 848 - 750 = 98.}
\end{center}
\hfill \textit{Arithmetic Progressions — General term from sum}
\vspace{0.3cm}

\question \textbf{26.} The 4th term of an AP is 0. Prove that the 25th term of the AP is three times its 11th term.

\textbf{Answer:}
\begin{center}
\fitmath{a_{25} = 3 \times a_{11}}
\end{center}

\textbf{Solution:}
\begin{enumerate}
  \item Express the 4th term as $a + 3d = 0$, leading to $a = -3d$.
  \item Express the 25th term as $a + 24d$.
  \item Express the 11th term as $a + 10d$.
  \item Substitute $a = -3d$ into both expressions and compare.
\end{enumerate}
\textbf{Derivation:}
\begin{center}
\fitmath{\begin{aligned}
a_4 = a + 3d = 0 \\
\implies a = -3d. \quad a_{25} = a + 24d = -3d + 24d = 21d. \quad a_{11} = a + 10d = -3d + 10d = 7d. \quad \text{Since } 21d = 3(7d), \text{ then } a_{25} = 3a_{11}.
\end{aligned}}
\end{center}
\hfill \textit{Arithmetic Progressions — Properties of AP}
\vspace{0.3cm}

\question \textbf{27.} How many terms of the AP 24, 21, 18, ... must be taken so that their sum is 78?

\textbf{Answer:}
\begin{center}
\fitmath{4 or 13}
\end{center}

\textbf{Solution:}
\begin{enumerate}
  \item Identify $a=24$, $d=-3$, $S_n=78$. Use the quadratic sum formula $S_n = n/2 [2a + (n-1)d]$.
  \item Substitute the known values to form a quadratic equation in $n$.
  \item Solve the quadratic equation for $n$ using factorization or the quadratic formula.
  \item Verify the two values of $n$.
\end{enumerate}
\textbf{Derivation:}
\begin{center}
\fitmath{\begin{aligned}
78 = \frac{n}{2} [2(24) + (n-1)(-3)] \\
\implies 156 = n(48 - 3n + 3) \\
\implies 156 = 51n - 3n^2 \\
\implies 3n^2 - 51n + 156 = 0 \\
\implies n^2 - 17n + 52 = 0 \\
\implies (n-4)(n-13) = 0. n=4, 13.
\end{aligned}}
\end{center}
\hfill \textit{Arithmetic Progressions — Sum of n terms}
\vspace{0.3cm}

\question \textbf{28.} Find the middle term of the Arithmetic Progression 7, 13, 19, ..., 241.

\textbf{Answer:}
\begin{center}
\fitmath{124}
\end{center}

\textbf{Solution:}
\begin{enumerate}
  \item Identify $a=7$, $d=6$, $l=241$. Find the total number of terms $n$ using $l = a + (n-1)d$.
  \item Determine if $n$ is odd or even to find the middle term position.
  \item Calculate the value of the middle term(s) using the formula $a_k = a + (k-1)d$.
\end{enumerate}
\textbf{Derivation:}
\begin{center}
\fitmath{\begin{aligned}
241 = 7 + (n-1)6 \\
\implies 234 = (n-1)6 \\
\implies n-1 = 39 \\
\implies n=40. \text{Since } n=40 \text{ is even, there are two middle terms: } 20^{th} \text{ and } 21^{st}. a_{20} = 7 + 19(6) = 121, a_{21} = 7 + 20(6) = 127. \text{Average is } 124.
\end{aligned}}
\end{center}
\hfill \textit{Arithmetic Progressions — General term}
\vspace{0.3cm}

\question \textbf{29.} The 4th term of an AP is 14 and the 12th term is 70. What is the first term $a$?

\textbf{Answer:}
(A)

\textbf{Solution:}
\begin{enumerate}
  \item Use the formula $a_n = a + (n-1)d$
  \item Set up two equations: $a + 3d = 14$ and $a + 11d = 70$
  \item Subtract the equations to find $d = 7$
  \item Substitute $d$ to find $a$
\end{enumerate}
\textbf{Derivation:}
\begin{center}
\fitmath{\begin{aligned}
a + 3d = 14 \\
a + 11d = 70 \\
(a + 11d) - (a + 3d) = 70 - 14 \\
8d = 56 \\
\implies d = 7 \\
a + 3(7) = 14 \\
\implies a = 14 - 21 = -7
\end{aligned}}
\end{center}
\hfill \textit{Arithmetic Progressions — nth term of an AP}
\vspace{0.3cm}

\question \textbf{30.} If the sum of first $n$ terms of an AP is $S_n = 3n^2 + n$, what is the common difference $d$?

\textbf{Answer:}
(C)

\textbf{Solution:}
\begin{enumerate}
  \item Find $S_1 = a_1 = 3(1)^2 + 1 = 4$
  \item Find $S_2 = a_1 + a_2 = 3(2)^2 + 2 = 14$
  \item Find $a_2 = S_2 - S_1 = 10$
  \item Find $d = a_2 - a_1 = 10 - 4 = 6$
\end{enumerate}
\textbf{Derivation:}
\begin{center}
\fitmath{\begin{aligned}
a_1 = S_1 = 4 \\
S_2 = 3(4) + 2 = 14 \\
a_2 = S_2 - S_1 = 14 - 4 = 10 \\
d = a_2 - a_1 = 10 - 4 = 6
\end{aligned}}
\end{center}
\hfill \textit{Arithmetic Progressions — Sum of n terms}
\vspace{0.3cm}

\question \textbf{31.} Which term of the AP: 21, 18, 15, ... is -81?

\textbf{Answer:}
(B)

\textbf{Solution:}
\begin{enumerate}
  \item Identify $a=21$ and $d=-3$
  \item Use $a_n = a + (n-1)d = -81$
  \item Solve for $n$
\end{enumerate}
\textbf{Derivation:}
\begin{center}
\fitmath{\begin{aligned}
21 + (n-1)(-3) = -81 \\
(n-1)(-3) = -102 \\
n-1 = 34 \\
n = 35
\end{aligned}}
\end{center}
\hfill \textit{Arithmetic Progressions — nth term of an AP}
\vspace{0.3cm}

\question \textbf{32.} In an AP, if $a=5$ and $d=3$, find the sum of the first 10 terms ($S_{10}$)?

\textbf{Answer:}
(C)

\textbf{Solution:}
\begin{enumerate}
  \item Use formula $S_n = \frac{n}{2}[2a + (n-1)d]$
  \item Substitute $n=10, a=5, d=3$
\end{enumerate}
\textbf{Derivation:}
\begin{center}
\fitmath{\begin{aligned}
S_{10} = \frac{10}{2}[2(5) + (10-1)3] \\
S_{10} = 5[10 + 27] \\
S_{10} = 5 \times 37 = 185
\end{aligned}}
\end{center}
\hfill \textit{Arithmetic Progressions — Sum of n terms}
\vspace{0.3cm}

\question \textbf{33.} If $k, 2k-1, 2k+1$ are three consecutive terms of an AP, find the value of $k$.

\textbf{Answer:}
(A)

\textbf{Solution:}
\begin{enumerate}
  \item Use property $b - a = c - b$ for three terms $a, b, c$ in AP
  \item Solve $(2k-1) - k = (2k+1) - (2k-1)$
\end{enumerate}
\textbf{Derivation:}
\begin{center}
\fitmath{\begin{aligned}
k-1 = 2 \\
k = 3
\end{aligned}}
\end{center}
\hfill \textit{Arithmetic Progressions — Properties of AP}
\vspace{0.3cm}

\question \textbf{34.} Find the sum of all odd numbers between 0 and 50.

\textbf{Answer:}
(B)

\textbf{Solution:}
\begin{enumerate}
  \item Identify AP: 1, 3, ..., 49
  \item Find $n$ using $a_n = a+(n-1)d$
  \item Use $S_n = \frac{n}{2}(a+l)$
\end{enumerate}
\textbf{Derivation:}
\begin{center}
\fitmath{\begin{aligned}
49 = 1 + (n-1)2 \\
\implies 48 = 2(n-1) \\
\implies n = 25 \\
S_{25} = \frac{25}{2}(1 + 49) = \frac{25}{2}(50) = 625
\end{aligned}}
\end{center}
\hfill \textit{Arithmetic Progressions — Sum of n terms}
\vspace{0.3cm}

\question \textbf{35.} The 10th term from the end of the AP: 5, 8, 11, ..., 185 is:

\textbf{Answer:}
(D)

\textbf{Solution:}
\begin{enumerate}
  \item Reverse the AP: 185, 182, ..., 5
  \item New $a=185, d=-3, n=10$
  \item Calculate $a_{10}$
\end{enumerate}
\textbf{Derivation:}
\begin{center}
\fitmath{\begin{aligned}
a_{10} = 185 + (10-1)(-3) \\
a_{10} = 185 - 27 = 158
\end{aligned}}
\end{center}
\hfill \textit{Arithmetic Progressions — nth term from the end}
\vspace{0.3cm}

\question \textbf{36.} If the sum of first $n$ terms is $n^2$, then the $n$th term is:

\textbf{Answer:}
(C)

\textbf{Solution:}
\begin{enumerate}
  \item $a_n = S_n - S_{n-1}$
  \item Substitute $S_n = n^2$ and $S_{n-1} = (n-1)^2$
\end{enumerate}
\textbf{Derivation:}
\begin{center}
\fitmath{\begin{aligned}
a_n = n^2 - (n-1)^2 \\
a_n = n^2 - (n^2 - 2n + 1) = 2n - 1
\end{aligned}}
\end{center}
\hfill \textit{Arithmetic Progressions — Sum of n terms}
\vspace{0.3cm}

\question \textbf{37.} How many two-digit numbers are divisible by 3?

\textbf{Answer:}
(A)

\textbf{Solution:}
\begin{enumerate}
  \item First number is 12, last is 99
  \item Use $a_n = a + (n-1)d$ where $d=3$
\end{enumerate}
\textbf{Derivation:}
\begin{center}
\fitmath{\begin{aligned}
99 = 12 + (n-1)3 \\
87 = (n-1)3 \\
29 = n-1 \\
\implies n = 30
\end{aligned}}
\end{center}
\hfill \textit{Arithmetic Progressions — Application of nth term}
\vspace{0.3cm}

\question \textbf{38.} What is the common difference of an AP in which $a_{25} - a_{20} = 45$?

\textbf{Answer:}
(B)

\textbf{Solution:}
\begin{enumerate}
  \item Express terms as $a+24d$ and $a+19d$
  \item Subtract the two expressions
\end{enumerate}
\textbf{Derivation:}
\begin{center}
\fitmath{\begin{aligned}
(a + 24d) - (a + 19d) = 45 \\
5d = 45 \\
\implies d = 9
\end{aligned}}
\end{center}
\hfill \textit{Arithmetic Progressions — nth term of an AP}
\vspace{0.3cm}

\question \textbf{39.} The sum of the first $n$ terms of an Arithmetic Progression is given by $S_n = 3n^2 + 5n$. Find the $n^{th}$ term of this AP and hence find the $20^{th}$ term.

\textbf{Answer:}
\begin{center}
\fitmath{6n+2; 122}
\end{center}

\textbf{Solution:}
\begin{enumerate}
  \item Use the relation $a_n = S_n - S_{n-1}$ to find the general term.
  \item Substitute $S_n = 3n^2 + 5n$.
  \item Calculate $S_{n-1} = 3(n-1)^2 + 5(n-1)$.
  \item Subtract $S_{n-1}$ from $S_n$ and simplify the expression.
  \item Substitute $n=20$ into the resulting formula to find $a_{20}$.
\end{enumerate}
\textbf{Derivation:}
\begin{center}
\fitmath{\begin{aligned}
a_n = S_n - S_{n-1} \\
a_n = (3n^2 + 5n) - [3(n-1)^2 + 5(n-1)] \\
a_n = 3n^2 + 5n - [3(n^2 - 2n + 1) + 5n - 5] \\
a_n = 3n^2 + 5n - [3n^2 - 6n + 3 + 5n - 5] \\
a_n = 3n^2 + 5n - 3n^2 + n + 2 = 6n + 2 \\
a_{20} = 6(20) + 2 = 122
\end{aligned}}
\end{center}
\hfill \textit{Arithmetic Progressions — Sum of n terms}
\vspace{0.3cm}

\question \textbf{40.} In an AP, the sum of first 10 terms is $-150$ and the sum of its next 10 terms is $-550$. Find the Arithmetic Progression.

\textbf{Answer:}
\begin{center}
\fitmath{a=4, d=-3; 4, 1, -2, -5, ...}
\end{center}

\textbf{Solution:}
\begin{enumerate}
  \item Set up the equation for $S_{10} = -150$.
  \item Recognize that the sum of the next 10 terms is $S_{20} - S_{10} = -550$.
  \item Find $S_{20} = -550 + (-150) = -700$.
  \item Solve the resulting system of two linear equations in $a$ and $d$.
  \item State the first few terms of the AP.
\end{enumerate}
\textbf{Derivation:}
\begin{center}
\fitmath{\begin{aligned}
S_{10} = 5(2a + 9d) = -150 \\
\implies 2a + 9d = -30 \quad (1) \\
S_{20} = 10(2a + 19d) = -700 \\
\implies 2a + 19d = -70 \quad (2) \\
(2) - (1): 10d = -40 \\
\implies d = -4 \text{ (Correction: d = -4, then 2a - 36 = -30, 2a = 6, a = 3)} \\
\text{Re-calculating: } 2a + 9(-4) = -30 \\
\implies 2a - 36 = -30 \\
\implies 2a = 6 \\
\implies a = 3 \\
\text{AP: } 3, -1, -5, -9, ...
\end{aligned}}
\end{center}
\hfill \textit{Arithmetic Progressions — Sum of n terms}
\vspace{0.3cm}

\end{questions}

\end{document}
